% GENERATED, edit the French unit or the translation table
\mxnewpage\section*{Solutions EV2-T1}\label{EV2-T1-CORR}
\textbf{MCQ:} 1 A, 2 D, 3 D, 4 B, 5 B, 6 C, 7 B, 8 C, 9 A, 10 C, 11 D, 12 C. Each correct answer: \num{0,5} point; total 6.
\paragraph{Checks of the fluency questions.} Order of operations; sum of fractions; subtraction of exponents; multiplying coefficient \num{1,2}; division by \num{0,8}; product \num{1,25} times \num{0,8}; solving an affine equation; vertical increment divided by horizontal; mean; angle divided by 360; complement to 1; substitution in brackets.
\subsection*{Exercise 1 -- Manufacturing and quality control}
\begin{enumerate}
\item First level: $\num{0,3}$ and $\num{0,7}$. After A: $\num{0,1}$ and $\num{0,9}$; after B: $\num{0,2}$ and $\num{0,8}$. Each pair sums to 1. \textit{Marking: 1 point. Result and basic justification expected; half a point possible for a correct but unfinished approach.}
\item $P(D)=\num{0,03}+\num{0,14}=\num{0,17}$; $P_D(A)=\dfrac{3}{17}$. \textit{Marking: 2 points. Method or justification: 1; calculation and conclusion: 1.}
\item No: $P_A(D)=\num{0,1}\ne\num{0,17}$. \textit{Marking: 2 points. Method or justification: 1; calculation and conclusion: 1.}
\end{enumerate}
\subsection*{Exercise 2 -- Minimum of a quotient}
\begin{enumerate}
\item $f\prime =\dfrac{(x-3)(x+3)}{x^2}$. \textit{Marking: 1 point. Result and basic justification expected; half a point possible for a correct but unfinished approach.}
\item Decreasing up to 3 then increasing; minimum 6. \textit{Marking: 2 points. Method or justification: 1; calculation and conclusion: 1.}
\item $1-\dfrac{9}{x^2}=\dfrac{3}{4}$ gives $x^2=36$, so x=6 only on the domain; point $\pt{6;\dfrac{15}{2}}$. \textit{Marking: 2 points. Method or justification: 1; calculation and conclusion: 1.}
\end{enumerate}
\subsection*{Exercise 3 -- Projection and circle}
\begin{enumerate}
\item $\pt{1;1}$. \textit{Marking: 1 point. Result and basic justification expected; half a point possible for a correct but unfinished approach.}
\item $\lambda =\dfrac{5+2-3}{2}=2$; H(3;0). \textit{Marking: 1 point. Result and basic justification expected; half a point possible for a correct but unfinished approach.}
\item $PH^2=8$; $(x-5)^2+(y-2)^2=8$. \textit{Marking: 2 points. Method or justification: 1; calculation and conclusion: 1.}
\end{enumerate}
\paragraph{Calibration.} Target duration: 20--25 min for the fluency questions, 80--85 min for the exercises, 10--15 min for proofreading. These durations are editorial estimates to be measured in class. Accept any correct method; do not penalise twice a numerical error that is propagated consistently. The proposed marking scheme is not a decision of the team to harmonise marking.
